\section{Discussion and Limitations}
\label{sec:discussion}

\paragraph{What the adaptation can establish.}
The intervention changes the input representation, not the pretrained 3D model.
Matched comparisons can test whether explicit line support helps a particular
backbone on reviewed patent crops. They cannot demonstrate a new multiview
conditioning architecture, hidden-state injection, or successful transfer to an
unevaluated dataset. The actual maps reveal a tradeoff: retaining thin contours
also retains dashed marks and stippled shading. These cues may be meaningful in
a patent but need not correspond to opaque 3D surfaces.

\paragraph{Measurement and coverage.}
Drawing-to-render metrics conflate structure with appearance and pose.
A high score may coexist with incorrect depth or disconnected geometry.
Thresholded ink masks are particularly different from filled rendered silhouettes.
We therefore withhold silhouette-based accuracy claims and separate complete-case
quality from generation and evaluation success. Patent-bootstrap intervals
characterize variation over sampled patents; they do not make a small convenience
sample representative, nor do they correct for selecting favorable parameters.

\paragraph{Reproducibility and pending evidence.}
The run manifests preserve crops, parameters, hashes, seeds, and status logs;
software and checkpoint revisions remain essential to a final comparison.
The pilot is implemented but has not yielded verified adaptation results in the
local project. Independent SketchDream and Know3D runs, held-out-patent evaluation,
and systematic failure selection remain necessary before reporting superiority.
Historical adjusted or simulated CSVs are excluded from the measured results.

\paragraph{Future evaluation.}
Further work should evaluate the four configurations on matched development
patents, examine distance-support and smoothing parameters, and freeze all
settings before a held-out patent test. Independent external-method runs and
multiple seeds are needed to assess relative performance and variability.
The configuration design and prepared inputs are not substitutes for those
measurements.

\section{Conclusion}
We formulate a controlled input-adaptation experiment for frozen TRELLIS and
analyze the patent reconstructions generated in this study. The method separates standard
preprocessing, background-removal bypass, distance-based line support, and
smoothing. Current evidence supports the archive characterization and evaluation
diagnostics; the pending matched experiments must determine whether adaptation
improves reconstruction agreement and at what computational or structural cost.
